\documentclass[11pt,a4paper]{article}
\usepackage[spanish]{babel}
\usepackage[T1]{fontenc}
\usepackage[margin=2cm,headheight=14pt,headsep=18pt,footskip=22pt]{geometry}
\usepackage{xcolor}
\usepackage{fancyhdr}
\usepackage{lastpage}
\usepackage{titlesec}
\usepackage{lmodern}
\usepackage[hidelinks]{hyperref}

\renewcommand{\familydefault}{\sfdefault}
\definecolor{VerdeUteq}{HTML}{006537}
\definecolor{AzulTitulo}{HTML}{19598A}
\definecolor{Dorado}{HTML}{D89D0B}

\newcommand{\estudiante}{Jonathan Holguín}
\newcommand{\asignatura}{Ingeniería de Requerimientos (ISR-401)}

\pagestyle{fancy}
\fancyhf{}
\fancyhead[L]{\small UTEQ \textperiodcentered\ \asignatura}
\fancyhead[R]{\small TA-IND-01 \textperiodcentered\ SPA 2026--2027}
\fancyfoot[L]{\small \estudiante}
\fancyfoot[R]{\small Página \thepage\ de \pageref{LastPage}}
\renewcommand{\footrulewidth}{0.4pt}

\titleformat{\section}
  {\large\bfseries}{\thesection.}{0.9em}{\color{VerdeUteq}}

\begin{document}

\begin{center}
  {\Large\bfseries\color{VerdeUteq} Universidad Técnica Estatal de Quevedo\par}
  \vspace{0.3em}
  {\small Facultad de Ciencias de la Computación
    \textcolor{Dorado}{\textbullet} Carrera de Software\par}
  \vspace{1em}
  {\large\bfseries\color{AzulTitulo} ACTIVIDAD SUMATIVA\par}
  \vspace{0.25em}
  {\large\bfseries TA-IND-01\par}
  {\Large\bfseries\color{VerdeUteq}
    Trabajo individual: IEEE 29148 y tipos de requisitos de un sistema conocido\par}
  \vspace{0.6em}
  {\normalsize \estudiante\par}
\end{center}
\noindent\textcolor{Dorado}{\rule{\linewidth}{0.8pt}}

\section{Sistema elegido}

GitHub Desktop para Windows permite trabajar con repositorios Git desde una
interfaz gráfica. Lo usan estudiantes y desarrolladores para revisar sus
cambios antes de guardarlos en un commit y compartirlos en GitHub%
~\cite{githubDesktop}. El análisis se centrará en tres acciones que forman
parte de un mismo flujo de trabajo: revisar qué cambió en los archivos,
organizar esos cambios en ramas y sincronizar el repositorio local con
GitHub.

\section{Resumen analítico de ISO/IEC/IEEE 29148:2018}

\section{Glosario}

\section{Requisitos y clasificación}

\section{Autoevaluación de los requisitos}

\section{Reflexión}

\section{Uso de IA}

\renewcommand{\refname}{Referencias}
\bibliographystyle{IEEEtran}
\bibliography{referencias}

\end{document}
